\documentclass[10pt,a4paper]{amsart}
\usepackage[margin=19mm]{geometry}
\usepackage{amsmath,amssymb,mathtools}
\usepackage[scr=rsfs,cal=euler]{mathalfa}
\usepackage{xfrac}
\usepackage[unicode,hidelinks,bookmarks=false]{hyperref}
\newcommand{\ie}{i.e.\ }
\newcommand{\cf}{cf.\ }
\newcommand{\resp}{respectively\ }
\newcommand{\Subsection}{\subsection}
\providecommand{\nobreakdash}{\nobreak\hbox{-}}
\setlength{\emergencystretch}{2em}
\multlinegap0pt
\usepackage{xfrac}
\usepackage{nicefrac}
\usepackage{amssymb}
\usepackage{tensor}
\usepackage{mathtools}
\usepackage{dsfont}
\usepackage[shortlabels]{enumitem}
\usepackage{bbm}
\newtheorem{prop}{Proposition}
\newcommand{\E}{\mathbb{E}}
\newcommand{\F}{\mathbb{F}}
\newcommand{\Gal}{\Gamma_l}
\newcommand{\Gau}{\Gamma_u}
\newcommand{\del}{\partial}
\newcommand{\dt}{\del_t}
\newcommand{\dz}{\del_z}
\newcommand{\eps}{\varepsilon}
\newcommand{\per}{\mathrm{per}}
\newcommand{\rB}{\mathrm{B}}
\newcommand{\rH}{\mathrm{H}}
\newcommand{\rL}{\mathrm{L}}
\newcommand{\rLp}{\rL^p}
\newcommand{\rLq}{\rL^q}
\newcommand{\rW}{\mathrm{W}}
\newcommand{\rX}{\mathrm{X}}
\newcommand{\srd}{\, \mathrm{d}}
\newcommand{\tand}{\enspace \text{and} \enspace}
\newcommand{\tin}{\enspace \text{in} \enspace}
\newcommand{\ton}{\enspace \text{on} \enspace}
\newcommand{\twith}{\enspace \text{with} \enspace}
\title{Moisture dynamics with phase changes coupled to heat-conducting, compressible fluids}
\author{Felix Brandt, Matthias Hieber, Lin Ma, Tarek Z\"ochling}
\date{Source: arXiv:2512.14931v1 (CC BY 4.0)}
\begin{document}
\maketitle

\noindent\emph{Source and licence.} Moisture dynamics with phase changes coupled to heat-conducting, compressible fluids. Version arXiv:2512.14931v1 (CC BY 4.0), distributed by arXiv under the Creative Commons Attribution 4.0 International licence, http://creativecommons.org/licenses/by/4.0/, as recorded in the arXiv licence metadata for this exact version and shown on the abstract page https://arxiv.org/abs/2512.14931v1. Full licence notice: \texttt{licenses/arxiv-2512.14931-CC-BY-4.0.txt}.\par\smallskip

\noindent\textit{Editorial scope.} Counted unit: the article's prop q (source0.tex lines 1532--1543). The article's complete original proof of it is delivered with it (source0.tex lines 1545--1602, one \texttt{proof} environment of 58 lines). No mathematics is written, repaired or re-derived: the statement and the proof are the article's own lines, in the article's own notation, and no retained line is substituted or re-flowed.  Retained uncounted as the source-local context the proof needs: lines 1515--1528.  Nothing was omitted from any retained span, so the package carries no omission record.  No retained line contains a citation, so the wrapper carries no bibliography and no bibliography entry.  No retained line contains a figure environment, an image-inclusion command or drawing code, so no image is delivered and the package declares no asset dependency.  Editorial formatting only, in addition: an AMS \texttt{amsart} preamble that restates the article's own declarations and macros used by the retained lines, namely the environments \texttt{prop} on separate counters, together with the standard packages those lines need.  The retained environments print in this document as Proposition 1 in order of appearance, whereas the article prints the same items with the numbering of its own full document.  The complete native source of the article as distributed in the arXiv e-print for this exact version is bundled unchanged as \texttt{sources/191131/source0.tex}.
\begin{equation}\label{eq: auxillary for q}
    \left\{
    \begin{aligned}
        \dt Q - \Delta Q
        &= f -Q^+, &&\tin (0,\tau) \times \Omega,\\
        \dz Q
        &= g_u, &&\ton (0,\tau) \times \Gau,\\
       \dz Q
        &= g_l, &&\ton (0,\tau) \times \Gal,\\
        Q(0)
        &= Q_0, &&\tin \Omega,
    \end{aligned}
    \right.
\end{equation}

\begin{prop}
    \label{prop: est on q}
    Let $\tau >0$ and $p,q \in [2,\infty)$. Assume that for $i \in \{ u,l\}$, the data satisfies
    \begin{equation}\label{eq: assu q}
        \begin{aligned}
            &f \in \E_{0,h,\tau}, \ g_i \in \F_{h,\tau}   \tand Q_0 \in \rB_{q,p,\per}^{2(1-\nicefrac{1}{p})}(\Omega) \twith\\
            &\| f \|_{\E_{0,h,\tau}} + \| g_i \|_{\F_{h,\tau}} +\|Q_0 \|_{\rB^{2(1-\nicefrac{1}{p})}_{q,p}(\Omega) } \leq C_0\eps^2 ,
        \end{aligned}
    \end{equation}
    for some $\eps>0$ and $C_0>0$. 
    Then there is a constant $C>0$ with $\| Q \|_{\E_{1,h,\tau}} \leq C \eps^2$.
\end{prop}

\begin{proof}
For brevity, for $q \in [1,\infty]$, we also simply denote the spatial $q$-norm by~$\| \cdot \|_q$.
We multiply \eqref{eq: auxillary for q} by $Q|Q|^{q-2}$ and integrate in space to obtain
\begin{equation*}
    \begin{aligned}
        \dt \| Q \|^q_q + \| \nabla | Q |^{\frac{q}{2}} \|^2_2 + \| |\nabla Q|  | Q|^{\frac{q-2}{2}} \|^2_2
        \leq C   \Bigl ( \int_\Omega \bigl (f - Q^+ \bigr )\cdot Q |Q|^{q-2} \srd x  + \int_G (g_u -g_l) \cdot Q |Q|^{q-2} \srd S\Bigr )
    \end{aligned}
\end{equation*}
for some constant $C>0$. 
In view of H\"older's and Young's inequality, we get
\begin{equation*}
    \Bigl | \int_\Omega (f - Q^+)\cdot Q |Q|^{q-2} \srd x \Bigr | \leq C \bigl ( \| f \|_q \cdot  \| Q \|_q^{q-1} + \| Q \|^q_q \bigr ) \leq C \bigl ( \| f\|^q_q + \| Q\|^q_q \bigr).
\end{equation*}
Concerning the boundary terms, we observe that similar arguments yield
\begin{equation*}
    \begin{aligned}
        \Bigl | \int_G (g_u -g_l) \cdot Q |Q|^{q-2} \srd S\Bigr | 
        \leq C \cdot \| g_u -g_l \|_{\rLq(G)}^q + \eps \cdot \| Q \|^q_{\rLq(G)}
        \leq C \cdot \| g_u -g_l \|^q_{\rB_{q,q}^{1 - \nicefrac{1}{q}}(G)} + \eps \cdot \| |Q|^{\frac{q}{2}} \|^2_{\rL^2(G)}.
    \end{aligned}
\end{equation*}
The trace theorem, interpolation and Young's inequality then imply
\begin{equation*}
    \begin{aligned}
        \| |Q|^{\frac{q}{2}} \|^2_{\rL^2(G)} \leq C \| |Q|^{\frac{q}{2}} \|^2_{\rW^{\nicefrac{1}{2},2}(\Omega)} 
        \leq C \bigl (\| |Q|^{\frac{q}{2}} \|^2_{\rH^{1}(\Omega)} + \||Q|^{\frac{q}{2}} \|^2_2 \bigr )
        \leq C\bigl ( \|\nabla |Q|^{\frac{q}{2}} \|^2_2 + \| Q\|^q_q \bigr ). 
    \end{aligned}
\end{equation*}
Hence, after absorbing, we obtain 
\begin{equation*}
    \dt \| Q \|^q_q \leq C \Bigl ( \| f\|^q_q + \| g_u -g_l \|^q_{\rB_{q,q}^{1 - \nicefrac{1}{q}}(G)}+ \| Q\|^q_q \Bigr). 
\end{equation*}
An integration in time then leads to
\begin{equation*}
    \| Q \|^q_q \leq C \bigl (\| f \|^q_{\rLq_\tau(\rL_x^q)} +\| g_u -g_l \|^q_{\rLq_\tau(\rB_{q,q}^{1 - \nicefrac{1}{q}}(G))}\bigr )+ C \int_0^t \| Q \|^q_q \ \mathrm{d}s + \| Q_0 \|^q_q.
\end{equation*}
By Gronwall's inequality, we conclude that
\begin{equation*}
    \| Q \|^q_q \leq \Bigl ( C  \| f \|^q_{\rLq_\tau(\rL_x^q)} +C\| g_u -g_l \|^q_{\rLq_\tau(\rB_{q,q}^{1 - \nicefrac{1}{q}}(G))} + \| Q_0 \|^q_q \Bigr ) \mathrm{e}^{C\tau}.
\end{equation*}
By Sobolev embeddings, we get $\rB^{2(1-\nicefrac{1}{p})}_{q,p,\per}(\Omega) \hookrightarrow \rLq(\Omega)$, so $\| Q\|_{\rL^\infty_\tau(\rL_x^q)} \leq C \eps^2$ by the assumptions on the data.
In particular, for all $q \geq 2$, we have 
\begin{equation*}
     \| \nabla | Q |^{\frac{q}{2}} \|^{\nicefrac{q}{2}}_2 + \| |\nabla Q|  | Q|^{\frac{q-2}{2}} \|^{\nicefrac{q}{2}}_2 \leq C \eps^2.
\end{equation*}
By the maximal $\rLp$-regularity on $\rX_0$, there exists a constant $C>0$ such that the local solution $Q \in \E_{1,h,\tau}$ to~\eqref{eq: auxillary for q} can be estimated by 
   \begin{equation*}
       \| Q \|_{\E_{1,h,\tau}} \leq C \bigl ( \| f \|_{\E_{0,h,\tau}} + \| Q^+\|_{\E_{0,h,\tau}}
       + \| g_u \|_{\F_{h,\tau}} +\| g_l \|_{\F_{h,\tau}} + \| Q_0 \|_{\rB_{q,p}^{2(1-\nicefrac{1}{p})}(\Omega)} \bigr ).
   \end{equation*}
The previous step joint with the embedding $\rL^\infty(0,\tau) \hookrightarrow \rLp(0,\tau)$ then yields
\begin{equation*}
    \| Q^+\|_{\E_{0,h,\tau}} \leq C \cdot \| Q \|_{\rL^\infty(0,\tau;\rLq(\Omega)) } \leq C\eps^2.
\end{equation*}
The claimed estimate then follows from the assumptions on the data.
\end{proof}

\end{document}
